\documentclass[12pt,a4paper,openany]{book}
\usepackage{geometry}
\usepackage{xepersian}
\usepackage{hyperref}
\usepackage{enumitem}
\usepackage{microtype}
\geometry{top=2.5cm,bottom=2.5cm,left=3cm,right=2.5cm}
\settextfont{Noto Naskh Arabic}
\setlatintextfont{Noto Sans}
\setdigitfont{Noto Naskh Arabic}
\hypersetup{colorlinks=true,linkcolor=black,urlcolor=blue}
\setlength{\parskip}{0.7em}
\setlength{\parindent}{0pt}

\title{مستندات جامع هوش مصنوعی و سیستم‌ها}
\author{AI Book Factory}
\date{\today}

\begin{document}
\maketitle
\tableofcontents
\clearpage

\chapter{مقدمه‌ای بر پردازش توزیع‌شده}

\section{مقدمه}

در دهه‌های گذشته، افزایش حجم داده‌ها و پیچیدگی محاسبات، مرزهای توانمندی سخت‌افزارهای تک‌سیستمی را جابه‌جا کرده است. قانون مور (Moore's Law) که سال‌ها پیش‌بینی‌کننده رشد دوبرابری تعداد ترانزیستورها در یک تراشه در هر دو سال بود، اکنون با محدودیت‌های فیزیکی و حرارتی مواجه شده است. در چنین شرایطی، برای آموزش مدل‌های هوش مصنوعی غول‌پیکر، تحلیل داده‌های کلان (Big Data) و میزبانی از سرویس‌های نرم‌افزاری با میلیون‌ها کاربر همزمان، تکیه بر یک ابررایانه تک‌تراشه‌ای دیگر پاسخگو نیست.

راه‌حل این چالش، «پردازش توزیع‌شده» (Distributed Processing) است. پردازش توزیع‌شده به معنای تقسیم یک کار محاسباتی بزرگ به بخش‌های کوچک‌تر و اجرای همزمان آن‌ها روی مجموعه‌ای از رایانه‌های مجزا است که از طریق شبکه با یکدیگر ارتباط برقرار می‌کنند. این رویکرد به سیستم‌ها اجازه می‌دهد تا بدون محدود شدن به سخت‌افزار یک ماشین واحد، به ظرفیت‌های محاسباتی و ذخیره‌سازی بی‌سابقه‌ای دست یابند.

در این فصل، به بررسی مفاهیم پایه‌ای، اصول معماری، چالش‌های طراحی و الگوهای رایج در سیستم‌های توزیع‌شده می‌پردازیم تا زیربنای لازم برای درک سیستم‌های پیچیده‌تر هوش مصنوعی و زیرساخت‌های مقیاس‌پذیر فراهم شود.

---

\section{۱. مفاهیم پایه‌ای پردازش توزیع‌شده}

برای ورود به دنیای سیستم‌های توزیع‌شده، ابتدا باید مفاهیم کلیدی و اصطلاحات پایه‌ای این حوزه را به دقت تعریف کنیم.

\subsection{پردازش موازی در برابر پردازش توزیع‌شده}

اگرچه در بسیاری از متون این دو اصطلاح به جای یکدیگر استفاده می‌شوند، اما تفاوت ساختاری مهمی بین آن‌ها وجود دارد: 	extit{ 	extbf{رایانش موازی (Parallel Computing):} معمولاً در یک سیستم واحد رخ می‌دهد که در آن چندین هسته پردازشی (CPU/GPU) به یک حافظه مشترک (Shared Memory) دسترسی دارند. زمان تاخیر ارتباطی در این سیستم‌ها بسیار ناچیز است. } 	extbf{پردازش توزیع‌شده (Distributed Processing):} شامل چندین سیستم مستقل (گره) است که هرکدام حافظه اختصاصی (Private Memory) خود را دارند و از طریق یک شبکه فیزیکی یا مجازی با یکدیگر پیام مبادله می‌کنند. در اینجا، تاخیر شبکه (Network Latency) یک فاکتور تعیین‌کننده است.

\subsection{گره (Node) و خوشه (Cluster)}

	extit{ 	extbf{گره (Node):} به هر واحد پردازشی مستقل در یک سیستم توزیع‌شده گفته می‌شود. یک گره می‌تواند یک سرور فیزیکی، یک ماشین مجازی (Virtual Machine) یا حتی یک کانتینر (Container) باشد. } 	extbf{خوشه (Cluster):} مجموعه‌ای از گره‌های متصل به هم که با هماهنگی یکدیگر کار می‌کنند تا به عنوان یک سیستم واحد و یکپارچه به نظر برسند.

\subsection{مقیاس‌پذیری (Scalability)}

یکی از اهداف اصلی پردازش توزیع‌شده، دستیابی به مقیاس‌پذیری است که به دو روش انجام می‌شود:

\begin{enumerate}
\item 	extbf{مقیاس‌پذیری عمودی (Vertical Scaling / Scale-up):} افزایش منابع سخت‌افزاری (مانند CPU، رم یا دیسک) در یک گره واحد. این روش محدودیت فیزیکی و هزینه تصاعدی دارد.
\item 	extbf{مقیاس‌پذیری افقی (Horizontal Scaling / Scale-out):} افزودن گره‌های بیشتر به خوشه. این رویکرد، هسته اصلی پردازش توزیع‌شده است و محدودیت‌های فیزیکی مقیاس‌دهی عمودی را ندارد.
\texttt{}\texttt{ مقیاس‌پذیری عمودی: [سرور کوچک] ───> [سرور بزرگتر] مقیاس‌پذیری افقی:   [سرور ۱] ───> [سرور ۱] + [سرور ۲] + [سرور ۳] }\texttt{}

\end{enumerate}
\subsection{تحمل خطا (Fault Tolerance) و افزونگی (Redundancy)}

در یک سیستم توزیع‌شده با هزاران گره، خرابی سخت‌افزار یک امر قطعی و روزمره است، نه یک استثنا. سیستم توزیع‌شده باید به گونه‌ای طراحی شود که خرابی یک یا چند گره منجر به توقف کل سیستم نشود. این هدف از طریق «افزونگی» (تکثیر داده‌ها و وظایف روی گره‌های مختلف) حاصل می‌شود.

\subsection{قضیه کپ (CAP Theorem)}

این قضیه که توسط «اریک برور» (Eric Brewer) مطرح شد، بیان می‌کند که یک سیستم توزیع‌شده ذخیره‌سازی داده همزمان فقط می‌تواند دو مورد از سه ویژگی زیر را تضمین کند: 	extit{ 	extbf{سازگاری (Consistency - C):} همه گره‌ها در یک لحظه داده‌های یکسانی را مشاهده کنند. } 	extbf{در دسترس بودن (Availability - A):} هر درخواستی که به سیستم ارسال می‌شود، بدون خطا پاسخ دریافت کند (حتی اگر داده‌ها به‌روز نباشند). * 	extbf{تحمل تقسیم شبکه (Partition Tolerance - P):} سیستم به کار خود ادامه دهد، حتی اگر ارتباط شبکه‌ای بین گره‌ها قطع یا دچار اختلال شدید شود.

از آنجا که در شبکه‌های واقعی احتمال بروز اختلال (Partition) همیشه وجود دارد، طراحان سیستم باید بین سازگاری و در دسترس بودن (انتخاب بین CP یا AP) یکی را برگزینند.

| نوع سیستم | اولویت | توصیف | نمونه کاربردی | | :--- | :--- | :--- | :--- | | 	extbf{CP} | سازگاری و تحمل تقسیم | در صورت بروز مشکل شبکه، سیستم پاسخ به درخواست‌ها را متوقف می‌کند تا از ناسازگاری داده جلوگیری کند. | تراکنش‌های بانکی | | 	extbf{AP} | در دسترس بودن و تحمل تقسیم | سیستم همیشه پاسخ می‌دهد، اما ممکن است برخی گره‌ها داده‌های قدیمی را برگردانند. | بخش نظرات شبکه‌های اجتماعی |

---

\section{۲. معماری و روند کار در سیستم‌های توزیع‌شده}

سیستم‌های توزیع‌شده بر اساس الگوهای معماری مشخصی طراحی می‌شوند تا هماهنگی بین گره‌ها بهینه‌سازی شود. در ادامه به دو معماری پرکاربرد و روندهای ارتباطی درون آن‌ها می‌پردازیم.

\subsection{الف) معماری ارباب-کارگر (Master-Worker / Leader-Follower)}

در این معماری، وظایف مدیریتی و اجرایی تفکیک می‌شوند: 	extit{ 	extbf{گره ارباب (Master):} وظیفه مدیریت وضعیت خوشه، تقسیم کارها بین گره‌های کارگر، پایش وضعیت سلامت (Health Check) گره‌ها و هماهنگی‌های کلی را بر عهده دارد. } 	extbf{گره‌های کارگر (Worker):} وظایف محاسباتی یا ذخیره‌سازی محول‌شده از سوی ارباب را دریافت، اجرا و نتیجه را بازمی‌گردانند.

این معماری در سیستم‌هایی مانند هادوف (Hadoop HDFS) و کوبرنتیز (Kubernetes) استفاده می‌شود. نقطه ضعف اصلی این الگو، احتمال تبدیل شدن گره ارباب به «نقطه تکی خرابی» (Single Point of Failure - SPOF) است که معمولاً با راه‌اندازی گره‌های پشتیبان (Standby Masters) حل می‌شود.

\subsection{ب) معماری همتا-به-همتا (Peer-to-Peer - P2P)}

در این الگو، هیچ گره متمایز یا کنترل‌کننده مرکزی وجود ندارد. همه گره‌ها نقش یکسانی ایفا می‌کنند و همزمان هم سرویس‌دهنده و هم سرویس‌گیرنده هستند. این معماری مقیاس‌پذیری فوق‌العاده بالا و تحمل خطای بی‌نظیری دارد، اما مدیریت همگام‌سازی و توافق در آن پیچیده‌تر است. شبکه‌های بلاک‌چین و سیستم‌های توزیع فایل مانند بیت‌تورنت از این الگو استفاده می‌کنند.

\subsection{روند ارتباطات و پیام‌رسانی}

گره‌ها در سیستم توزیع‌شده باید بتوانند با یکدیگر ارتباط برقرار کنند. دو روش اصلی برای این کار وجود دارد:

\begin{enumerate}
\item 	extbf{فراخوانی از راه دور (Remote Procedure Call - RPC):} به برنامه‌نویس اجازه می‌دهد متدی را روی یک ماشین دیگر اجرا کند، به گونه‌ای که گویا در حال اجرای آن روی سیستم محلی خود است. پروتکل gRPC یک نمونه مدرن و محبوب از این فناوری است.
\item 	extbf{صف‌های پیام (Message Queues):} ارتباط ناهمگام (Asynchronous) را از طریق واسط‌هایی مانند RabbitMQ یا Apache Kafka میسر می‌سازند. در این روش، فرستنده پیام خود را در صف قرار می‌دهد و نیازی به منتظر ماندن برای پاسخ فوری گیرنده ندارد. این الگو باعث کاهش وابستگی مستقیم (Decoupling) سیستم‌ها می‌شود.
\texttt{}\texttt{ [گره فرستنده] ───(ارسال پیام ناهمگام)───> [ صف پیام ] ───(دریافت پیام)───> [گره گیرنده] }\texttt{}

\end{enumerate}
\subsection{الگوریتم‌های اجماع (Consensus Algorithms)}

یکی از چالش‌های بزرگ در سیستم‌های بدون مرکزیت یا با پویایی بالا، توافق گره‌ها بر سر یک واقعیت واحد (مثلاً ترتیب تراکنش‌ها یا انتخاب رهبر جدید) است. برای حل این مشکل از الگوریتم‌های اجماع مانند 	extbf{Raft} یا 	extbf{Paxos} استفاده می‌شود. این الگوریتم‌ها تضمین می‌کنند که حتی در صورت خرابی بخشی از گره‌ها، کل سیستم روی وضعیت نهایی به توافق برسد.

---

\section{مثال عملی: شبیه‌سازی الگوی نقشه‌کاهش (MapReduce) در پایتون}

یکی از معروف‌ترین الگوهای برنامه‌نویسی و پردازش توزیع‌شده، الگوی 	extbf{نقشه‌کاهش (MapReduce)} است که ابتدا توسط گوگل معرفی شد. این الگو عملیات را به سه مرحله اصلی تقسیم می‌کند:

\begin{enumerate}
\item 	extbf{نگاشت (Map):} داده‌های ورودی بزرگ به تکه‌های کوچک‌تر تقسیم شده و هر تکه توسط یک گره پردازش می‌شود تا جفت‌های کلید-مقدار (Key-Value) تولید شوند.
\item 	extbf{درهم‌آمیزی (Shuffle):} داده‌های تولید شده بر اساس کلیدها گروه‌بندی و مرتب می‌شوند تا به گره‌های کاهنده مناسب ارسال شوند.
\item 	extbf{کاهش (Reduce):} مقادیر مربوط به هر کلید با یکدیگر ترکیب یا تجمیع می‌شوند تا خروجی نهایی به دست آید.
در ادامه، یک شبیه‌سازی ساده از این الگو را با استفاده از کتابخانه استاندارد \texttt{multiprocessing} در پایتون پیاده‌سازی می‌کنیم. این مثال فرآیند شمارش کلمات (Word Count) را در یک متن بزرگ به صورت موازی/توزیع‌شده شبیه‌سازی می‌کند.

\texttt{}`python from multiprocessing import Pool from collections import defaultdict import re

def clean\_text(text): """پاک‌سازی متن و استخراج کلمات.""" return re.findall(r'\textbackslash{}w+', text.lower())

def map\_worker(text\_chunk): """ مرحله نگاشت (Map): یک تکه متن را دریافت کرده و جفت‌های (word, 1) را تولید می‌کند. """ word\_counts = [] words = clean\_text(text\_chunk) for word in words: word\_counts.append((word, 1)) return word\_counts

def shuffle\_and\_group(mapped\_data): """ مرحله درهم‌آمیزی (Shuffle): داده‌های حاصل از مرحله نگاشت را بر اساس کلید (کلمه) گروه‌بندی می‌کند. """ grouped\_data = defaultdict(list) for sublist in mapped\_data: for word, count in sublist: grouped\_data[word].append(count) return grouped\_data.items()

def reduce\_worker(item): """ مرحله کاهش (Reduce): لیست مقادیر مربوط به یک کلید را دریافت کرده و آن‌ها را جمع می‌زند. """ word, counts = item return word, sum(counts)

if \_\_name\_\_ == "\_\_main\_\_":

\end{enumerate}
\chapter{متن ورودی شبیه‌سازی شده (داده‌های بزرگ)}

documents = [ "Distributed systems are complex but highly scalable.", "Parallel processing is different from distributed processing.", "Scalability and fault tolerance are key concepts in distributed computing.", "Python multiprocessing can simulate distributed environments on a single machine." ]

print("--- شروع فرآیند پردازش توزیع‌شده (شبیه‌سازی MapReduce) ---")

\chapter{۱. مرحله نگاشت موازی با استفاده از Pool پردازنده‌ها (شبیه‌سازی گره‌های کارگر)}

with Pool(processes=4) as pool: mapped\_results = pool.map(map\_worker, documents)

print(f"\textbackslash{}n[مرحله Map] خروجی گره‌های کارگر (نمونه): \{mapped\_results[0][:3]\}...")

\chapter{۲. مرحله درهم‌آمیزی (Shuffle) برای گروه‌بندی کلمات مشابه}

shuffled\_results = shuffle\_and\_group(mapped\_results) print(f"\textbackslash{}n[مرحله Shuffle] داده‌ها بر اساس کلیدها گروه‌بندی شدند.")

\chapter{۳. مرحله کاهش موازی (Reduce) برای محاسبه جمع نهایی هر کلمه}

with Pool(processes=4) as pool: final\_counts = pool.map(reduce\_worker, shuffled\_results)

\chapter{نمایش نتایج نهایی (مرتب‌شده بر اساس بیشترین تکرار)}

sorted\_results = sorted(final\_counts, key=lambda x: x[1], reverse=True)

print("\textbackslash{}n--- نتایج نهایی شمارش کلمات (تعدادی از کلمات پرتکرار) ---") for word, count in sorted\_results[:5]: print(f"کلمه '\{word\}': \{count\} بار") \texttt{}`

\subsection{تحلیل کد:}

	extit{ تابع \texttt{map\_worker} شبیه‌ساز کاری است که روی گره‌های مجزا انجام می‌شود. هر گره ورودی خود را مستقل از دیگران پردازش می‌کند. } کلاس \texttt{Pool} از کتابخانه \texttt{multiprocessing} پایتون، وظیفه مدیریت فرآیندها (Processes) را به عنوان گره‌های مجازی بر عهده دارد. * در مرحله نهایی، سیستم نتایج پراکنده را یکپارچه کرده و خروجی منسجم ارائه می‌دهد.

---

\section{خطاهای رایج در سیستم‌های توزیع‌شده}

طراحی و توسعه سیستم‌های توزیع‌شده با چالش‌های منحصر‌به‌فردی همراه است. در ادامه به برخی از رایج‌ترین خطاها و اشتباهات محاسباتی در این حوزه اشاره می‌کنیم:

\subsection{۱. فرضیات نادرست درباره شبکه (The Fallacies of Distributed Computing)}

بسیاری از توسعه‌دهندگان در ابتدای کار فرضیاتی را در ذهن خود دارند که در عمل نادرست هستند. این فرضیات به «مغالطه‌های رایانش توزیع‌شده» معروفند: 	extit{ 	extbf{فرض بر این که شبکه همیشه قابل اعتماد است:} قطعی موقت کابل‌ها، تداخل‌های رادیویی یا خطاهای سوئیچ‌ها همواره رخ می‌دهند. } 	extbf{فرض صفر بودن تاخیر (Latency):} ارسال داده روی شبکه زمان‌بر است. طراحی سیستمی که نیاز به پیام‌رسانی‌های مکرر و ریزدانه (Chatty API) دارد، بازدهی بسیار پایینی خواهد داشت. * 	extbf{فرض بی‌نهایت بودن پهنای باند:} پهنای باند شبکه محدود است و انتقال حجم عظیمی از داده‌ها بین گره‌ها می‌تواند به سرعت به گلوگاه (Bottleneck) سیستم تبدیل شود.

\subsection{۲. پدیده دوپارگی مغز (Split-Brain)}

در معماری‌های دارای گره رهبر (Leader)، اگر ارتباط شبکه‌ای میان گره‌ها قطع شود، ممکن است خوشه به دو بخش مجزا تقسیم شود. در این حالت، هر بخش ممکن است تصور کند رهبر اصلی از دست رفته است و اقدام به انتخاب یک رهبر جدید کند. وجود دو رهبر همزمان در یک سیستم توزیع‌شده منجر به بازنویسی داده‌های متناقض و تخریب شدید پایگاه داده می‌شود. * 	extbf{راه‌حل:} استفاده از حد نصاب یا کواروم (Quorum)؛ به این معنی که یک رهبر تنها زمانی رسمیت دارد که تایید اکثریت مطلق گره‌ها (مثلاً بیش از نصف) را داشته باشد.

\subsection{۳. شرایط رقابتی (Race Conditions) و عدم همگام‌سازی زمان}

در سیستم‌های توزیع‌شده، مفهوم «زمان مطلق» وجود ندارد. ساعت‌های داخلی سخت‌افزارها (حتی با استفاده از NTP) تفاوت‌های میلی‌ثانیه‌ای با یکدیگر دارند. اگر سیستم برای تعیین ترتیب رویدادها صرفاً به زمان محلی گره‌ها اتکا کند، خطاهای شدیدی در ثبت داده‌ها رخ می‌دهد. * 	extbf{راه‌حل:} استفاده از ساعت‌های منطقی (Logical Clocks) مانند ساعت لمپورت (Lamport Timestamps) یا ساعت‌های برداری (Vector Clocks) به جای ساعت‌های فیزیکی.

\subsection{۴. مهندسی بیش از حد (Over-engineering)}

استفاده از سیستم‌های توزیع‌شده پیچیدگی‌های عملیاتی، هزینه‌های نگهداری و دیباگ کردن را به شدت افزایش می‌دهد. پیاده‌سازی یک معماری توزیع‌شده برای پروژه‌هایی که داده‌های آن‌ها روی یک سرور مدرن واحد قابل پردازش است، یک خطای استراتژیک رایج است. همواره باید مطمئن شد که نیازهای پروژه فراتر از توان پردازش عمودی (Scale-up) است، سپس به سراغ سیستم‌های توزیع‌شده رفت.

---

\section{جمع‌بندی}

پردازش توزیع‌شده سنگ بنای زیرساخت‌های مدرن فناوری اطلاعات، سیستم‌های کلان‌داده و هوش مصنوعی مقیاس‌پذیر است. با گذر از محدودیت‌های فیزیکی تک‌تراشه‌ها، خوشه‌های توزیع‌شده به ما امکان می‌دهند تا محاسباتی در مقیاس پتابایت را انجام دهیم. با این حال، این قدرت بالا با چالش‌های بزرگی چون تاخیر شبکه، لزوم مدیریت تحمل خطا، هماهنگی گره‌ها و رعایت محدودیت‌های قضیه CAP همراه است.

درک دقیق مفاهیمی چون مقیاس‌پذیری افقی، الگوهای معماری مانند ارباب-کارگر و الگوریتم‌های اجماع، پیش‌نیاز ورود به مباحث پیشرفته‌تر نظیر آموزش توزیع‌شده مدل‌های یادگیری عمیق و طراحی خط‌لوله‌های داده (Data Pipelines) است که در فصول بعدی کتاب به آن‌ها خواهیم پرداخت.


\chapter{معماری گیت‌هاب اکشنز در اتوماسیون}

\chapter{فصل دوم: معماری گیت‌هاب اکشنز در اتوماسیون}

\section{مقدمه}

در فصل پیشین با اصول پردازش توزیع‌شده، نحوه مدیریت گره‌ها و الگوهای معماری برای مقیاس‌پذیری محاسبات آشنا شدیم. اما در دنیای واقعی مهندسی نرم‌افزار و هوش مصنوعی، مدیریت کدها، تست‌ها، آماده‌سازی داده‌ها و استقرار مدل‌ها نیازمند ابزاری است که بتواند فرایندهای پیچیده را به صورت خودکار و توزیع‌شده اجرا کند. اینجاست که ابزارهای «ادغام مداوم و تحویل مداوم» (Continuous Integration / Continuous Deployment - CI/CD) و پلتفرم‌های اتوماسیون وارد میدان می‌شوند.

یکی از قدرتمندترین و پرکاربردترین بسترهای اتوماسیون مدرن، 	extbf{گیت‌هاب اکشنز (GitHub Actions)} است. گیت‌هاب اکشنز صرفاً ابزاری برای اجرای تست‌های واحد نیست؛ بلکه یک سیستم توزیع‌شده مبتنی بر رویداد (Event-driven Distributed System) است که به توسعه‌دهندگان اجازه می‌دهد تمامی مراحل چرخه حیات نرم‌افزار را از زمان ثبت کد تا استقرار نهایی در محیط‌های ابری اتوماتیک کنند.

در این فصل، معماری داخلی گیت‌هاب اکشنز، اجزای کلیدی، نحوه توزیع و مدیریت کارها روی مجری‌ها (Runners)، مکانیزم‌های امنیتی و نحوه ساخت یک خط‌لوله (Pipeline) اتوماسیون پیشرفته بررسی می‌شود.

---

\section{۱. مفاهیم و مولفه‌های کلیدی گیت‌هاب اکشنز}

برای درک معماری گیت‌هاب اکشنز، ابتدا باید ساختار مفهومی و اصطلاحات پایه‌ای آن را بشناسیم. این بستر از اجزای متصلی تشکیل شده که نحوه تعامل آن‌ها در شکل زیر خلاصه شده است:

\texttt{}\texttt{ [رویداد (Event)] │ ▼ [گردش‌کار (Workflow)] │ ├───> [کار ۱ (Job 1)] ───> [مرحله ۱] ───> [اکشن / دستور] │                          [مرحله ۲] │ └───> [کار ۲ (Job 2)] ───> [مرحله ۱] }\texttt{}

\subsection{گردش‌کار (Workflow)}

یک روال خودکار قابل پیکربندی است که از یک یا چند «کار» تشکیل شده و به صورت یک فایل YAML در مسیر \texttt{.github/workflows} مخزن ذخیره می‌شود. گردش‌کار می‌تواند بر اساس رویدادهای مخزن، زمان‌بندی‌های مشخص یا به صورت دستی تحریک (Trigger) شود.

\subsection{رویداد (Event)}

هر فعالیت مشخص در مخزن گیت‌هاب که می‌تواند باعث اجرای یک گردش‌کار شود. رویدادها طیف وسیعی را شامل می‌شوند؛ مانند: 	extit{ ارسال کد جدید (\texttt{push}) } ایجاد درخواست کشش (\texttt{pull\_request}) 	extit{ انتشار یک نسخه جدید (\texttt{release}) } زمان‌بندی مبتنی بر کران‌جاب (\texttt{schedule}) * فراخوانی دستی یا از طریق API (\texttt{workflow\_dispatch})

\subsection{کار (Job)}

مجموعه‌ای از «مراحل» (Steps) است که روی یک مجری (Runner) واحد اجرا می‌شود. به‌طور پیش‌فرض، کارهای موجود در یک گردش‌کار به صورت 	extbf{موازی} اجرا می‌شوند، مگر اینکه بین آن‌ها وابستگی (Dependency) تعریف شده باشد. این ساختار موازی‌سازی مستقیم از مفاهیم پردازش توزیع‌شده الهام گرفته شده است.

\subsection{مرحله (Step)}

یک دستور اجرایی مستقل درون یک کار است. یک مرحله می‌تواند یک دستور شل (Shell Command) یا یک «اکشن» آماده باشد. همه مراحلِ یک کار روی یک ماشین واحد و در یک فضای کاری مشترک اجرا می‌شوند؛ بنابراین داده‌ها می‌توانند بین مراحل یک کار به راحتی به اشتراک گذاشته شوند.

\subsection{اکشن (Action)}

کوچک‌ترین بلوک ساختاری در گیت‌هاب اکشنز است. اکشن‌ها برنامه‌های پیش‌ساخته و قابل مجدداًستفاده‌ای هستند که وظایف خاصی (مانند بارگیری کد از مخزن، تنظیم محیط پایتون، یا ارسال پیام به اسلک) را انجام می‌دهند. اکشن‌ها می‌توانند از «بازارچه گیت‌هاب» (GitHub Marketplace) فراخوانی شوند یا درون خود مخزن نوشته شوند.

\subsection{مجری (Runner)}

سرور یا ماشینی است که برنامه «مجری گیت‌هاب» روی آن نصب شده و کارهای تعریف‌شده در گردش‌کار را اجرا می‌کند. مجری‌ها به دو دسته اصلی تقسیم می‌شوند:

\begin{enumerate}
\item 	extbf{مجری‌های میزبانی‌شده توسط گیت‌هاب (GitHub-hosted Runners):} ماشین‌های مجازی مدیریت‌شده توسط گیت‌هاب با سیستم‌عامل‌های لینوکس، ویندو یا مک‌او‌اس که پس از پایان هر کار کاملاً پاکسازی و نابود می‌شوند (Ephemeral).
\item 	extbf{مجری‌های اختصاصی (Self-hosted Runners):} ماشین‌های فیزیکی، مجازی یا کانتینرهایی که توسط خود کاربر در زیرساخت شخصی یا ابری مدیریت می‌شوند و گیت‌هاب اکشنز وظایف را به آن‌ها محول می‌کند.
| ویژگی | مجری میزبانی‌شده (GitHub-hosted) | مجری اختصاصی (Self-hosted) | | :--- | :--- | :--- | | 	extbf{مدیریت زیرساخت} | کامل بر عهده گیت‌هاب | بر عهده کاربر | | 	extbf{ایزوله‌سازی} | ماشین مجازی پاک‌سازی‌شونده برای هر کار | محیط پایدار (نیازمند پاکسازی دستی یا کانتینری) | | 	extbf{دسترسی به شبکه داخلی} | محدود (نیازمند VPN/IP Allowlist) | مستقیم به شبکه داخلی شرکت | | 	extbf{سخت‌افزار سفارشی (GPU/RAM بالا)} | محدود به پلن‌های خاص | کاملاً سفارشی و بدون محدودیت |

---

\end{enumerate}
\section{۲. معماری و روند کار در گیت‌هاب اکشنز}

معماری گیت‌هاب اکشنز بر پایه الگوی تفکیک «سطح کنترل» (Control Plane) و «سطح اجرا» (Execution Plane) بنا شده است. این طراحی سیستم را قادر می‌سازد تا میلیون‌ها گردش‌کار همزمان را بدون تداخل و با امنیت بالا مدیریت کند.

\texttt{}\texttt{ +-------------------------------------------------------------------+ |                       سطح کنترل (GitHub)                           | |  [ Event Service ] ──> [ Orchestrator Engine ] ──> [ Job Queue ]  | +-------------------------------------------------------------------+ ▲ Long-Polling / TLS ▼ +-------------------------------------------------------------------+ |                        سطح اجرا (Runner)                          | |  [ Runner Daemon ] ──> [ Worker Process ] ──> [ Steps / Actions ] | +-------------------------------------------------------------------+ }\texttt{}

\subsection{الف) سطح کنترل (Control Plane)}

سطح کنترل درون زیرساخت ابری گیت‌هاب قرار دارد و مسئولیت‌های زیر را بر عهده دارد: 	extit{ 	extbf{دریافت و تحلیل رویدادها (Event Ingestion):} هنگامی که اتفاقی در مخزن رخ می‌دهد (مانند Push)، سرویس رویداد آن را دریافت کرده و فایل‌های YAML مربوطه را بررسی می‌کند. } 	extbf{ارزیابی و زمان‌بندی (Orchestration \& Scheduling):} موتور ارکستراسیون شرایط اجرای گردش‌کار را می‌سنجد، متغیرها و ماتریس‌ها را ارزیابی کرده و کارها (Jobs) را تفکیک می‌کند. * 	extbf{صف کارهای در انتظار (Job Queueing):} کارهای آماده به اجرا در یک صف توزیع‌شده قرار می‌گیرند تا توسط مجری‌های مناسب دریافت شوند.

\subsection{ب) سطح اجرا (Execution Plane)}

سطح اجرا شامل ماشین‌هایی است که فرایند واقعی پردازش را انجام می‌دهند. ارتباط بین سطح کنترل و سطح اجرا بر اساس الگوی 	extbf{Pull} انجام می‌شود، نه Push.

	extbf{چرا الگوی Pull؟} مجری‌ها به صورت مداوم (از طریق Long-Polling یا WebSockets امن بر روی HTTPS) صف گیت‌هاب را بررسی می‌کنند تا کار جدیدی را دریافت کنند. این رویکرد دو مزیت کلیدی دارد:

\begin{enumerate}
\item 	extbf{امنیت شبکه:} نیازی به باز کردن پورت‌های ورودی (Inbound Ports) روی جدار firewall سرورهای مجری نیست.
\item 	extbf{مقیاس‌پذیری:} مجری‌های جدید می‌توانند به راحتی بدون نیاز به ثبت IP در سرور مرکزی، به خوشه متصل شده و شروع به کار کنند.
\end{enumerate}
\subsection{ج) چرخه حیات اجرای یک کار (Job Execution Lifecycle)}

\begin{enumerate}
\item 	extbf{تحریک رویداد (Triggering):} رویداد صادر شده و سرویس مرکزی گیت‌هاب آن را به یک یا چند کار تبدیل می‌کند.
\item 	extbf{تخصیص مجری (Runner Assignment):} سرویس مرکزی بر اساس برچسب‌ها (\texttt{runs-on: ubuntu-latest}) یک مجری آزاد را انتخاب یا یک ماشین مجازی جدید را آماده می‌کند.
\item 	extbf{بارگیری و آماده‌سازی (Environment Setup):} سرویس مجری (Runner Daemon) کار را دریافت کرده، متغیرهای محیطی و رازها (Secrets) را تزریق می‌کند و فضای کاری (Workspace) را آماده می‌سازد.
\item 	extbf{اجرای مراحل (Step Execution):} مراحل به ترتیب اجرا می‌شوند. اگر یک مرحله با خطا مواجه شود (مگر اینکه \texttt{continue-on-error} تنظیم شده باشد)، اجرای بقیه مراحل متوقف شده و کار به عنوان شکست‌خورده علامت‌گذاری می‌شود.
\item 	extbf{جمع‌آوری خروجی و پاکسازی (Post-processing \& Cleanup):} لاگ‌ها به صورت زنده به گیت‌هاب ارسال می‌شوند، مصنوعات (Artifacts) آپلود شده و محیط مجری پاکسازی می‌شود.
\end{enumerate}
\subsection{د) مدیریت رازها و امنیت (Secrets and Security)}

امنیت در سیستم‌های اتوماسیون توزیع‌شده اهمیت حیاتی دارد. گیت‌هاب اکشنز چند مکانیزم امنیتی اصلی را پیاده‌سازی می‌کند:

	extit{ 	extbf{ماسکه کردن رازها (Secret Masking):} تمامی متغیرهای حساسی که در بخش \texttt{Secrets} تنظیم می‌شوند، در زمان چاپ در لاگ‌ها به صورت اتوماتیک با \texttt{***} جایگزین می‌شوند. } 	extbf{توکن اختصاصی (\texttt{GITHUB\_TOKEN}):} در ابتدای هر کار، گیت‌هاب یک توکن احراز هویت موقت تولید می‌کند که تنها به همین کار اختصاص دارد و با پایان یافتن کار منقضی می‌شود. دسترسی‌های این توکن قابل محدودسازی (Fine-grained permissions) است. * 	extbf{پروتکل OIDC (OpenID Connect):} برای ارتباط با ابرهایی مانند AWS، Azure یا GCP، نیازی به ذخیره‌سازی کلیدهای دسترسی بلندمدت نیست. گیت‌هاب اکشنز از OIDC پشتیبانی می‌کند تا مجری بتواند توکن‌های کوتاه مدّت و امن را مستقیماً از ارائه‌دهنده ابر دریافت کند.

\subsection{هـ) مدیریت کش و مصنوعات (Caching and Artifacts)}

در خطوط لوله اتوماسیون، دانلود مجدد وابستگی‌ها (Dependencies) در هر بار اجرا، موجب اتلاف وقت و پهنای باند می‌شود. 	extit{ 	extbf{کش (Cache):} گیت‌هاب امکان ذخیره فایل‌های وابستگی (مانند پکیج‌های \texttt{pip} یا \texttt{npm}) را بر اساس کلیدهای هش فایلهای قفل (Lock files) فراهم می‌کند. } 	extbf{مصنوعات (Artifacts):} فایل‌ها یا داده‌های خروجی یک کار (مانند فایل‌های binary تولیدشده یا گزارش‌های تست) را می‌توان به عنوان Artifact ذخیره کرد تا کارهای بعدی یا کاربران به آن‌ها دسترسی داشته باشند.

---

\section{مثال عملی}

در این بخش، یک فایل کاملی از گردش‌کار گیت‌هاب اکشنز را بررسی می‌کنیم. این سناریو شامل اجرای موازی تست‌ها روی چند نسخه پایتون (Matrix Strategy)، استفاده از سیستم کش، ذخیره پوشش تست به عنوان Artifact و در نهایت ساخت تصویر داکر است.

مسیر فایل: \texttt{.github/workflows/ci-cd-pipeline.yml}

\texttt{}`yaml name: Production CI/CD Pipeline

\chapter{تعیین رویدادهای تحریک‌کننده}

on: push: branches: [ "main" ] pull\_request: branches: [ "main" ]

\chapter{تنظیم دسترسی‌های توکن امنیتی}

permissions: contents: read packages: write

jobs:

\chapter{کار اول: اجرای تست‌ها به صورت ماتریسی}

test: name: Run Unit Tests (Python \$\{\{ matrix.python-version \}\}) runs-on: ubuntu-latest

strategy: matrix: python-version: ["3.10", "3.11"]

steps:

\chapter{۱. دریافت کد از مخزن}

\begin{itemize}
\item name: Checkout Code
uses: actions/checkout@v4

\end{itemize}
\chapter{۲. تنظیم محیط پایتون}

\begin{itemize}
\item name: Set up Python
uses: actions/setup-python@v5 with: python-version: \$\{\{ matrix.python-version \}\}

\end{itemize}
\chapter{۳. بازگردانی یا ایجاد حافظه نهان (Cache) برای وابستگی‌ها}

\begin{itemize}
\item name: Cache Python Dependencies
uses: actions/cache@v4 with: path: \textasciitilde{}/.cache/pip key: \$\{\{ runner.os \}\}-pip-\$\{\{ hashFiles('**/requirements.txt') \}\} restore-keys: | \$\{\{ runner.os \}\}-pip-

\end{itemize}
\chapter{۴. نصب وابستگی‌ها}

\begin{itemize}
\item name: Install Dependencies
run: | python -m pip install --upgrade pip if [ -f requirements.txt ]; then pip install -r requirements.txt; fi pip install pytest pytest-cov

\end{itemize}
\chapter{۵. اجرای تست‌ها و تولید گزارش}

\begin{itemize}
\item name: Run Tests with Pytest
run: | pytest --cov=src --cov-report=xml:coverage.xml

\end{itemize}
\chapter{۶. آپلود گزارش پوشش تست به عنوان Artifact (تنها برای یک نسخه)}

\begin{itemize}
\item name: Upload Test Coverage Artifact
if: matrix.python-version == '3.11' uses: actions/upload-artifact@v4 with: name: code-coverage-report path: coverage.xml

\end{itemize}
\chapter{کار دوم: ساخت و ارسال تصویر داکر (تنها پس از موفقیت کار تست)}

build-and-publish: name: Build \& Push Docker Image needs: test if: github.ref == 'refs/heads/main' \&\& github.event\_name == 'push' runs-on: ubuntu-latest

steps:

\begin{itemize}
\item name: Checkout Code
uses: actions/checkout@v4

\end{itemize}
\chapter{تنظیم داکر بیلداوت (Buildx)}

\begin{itemize}
\item name: Set up Docker Buildx
uses: docker/setup-buildx-action@v3

\end{itemize}
\chapter{ورود به Container Registry گیت‌هاب}

\begin{itemize}
\item name: Log in to GitHub Container Registry
uses: docker/login-action@v3 with: registry: ghcr.io username: \$\{\{ github.actor \}\} password: \$\{\{ secrets.GITHUB\_TOKEN \}\}

\end{itemize}
\chapter{ساخت و ارسال تصویر}

\begin{itemize}
\item name: Build and Push Docker Image
uses: docker/build-push-action@v5 with: context: . push: true tags: | ghcr.io/\$\{\{ github.repository \}\}/app:latest ghcr.io/\$\{\{ github.repository \}\}/app:\$\{\{ github.sha \}\} cache-from: type=gha cache-to: type=gha,mode=max \texttt{}`

\end{itemize}
\subsection{تحلیل جریان اجرای مثال فوق:}

\begin{enumerate}
\item 	extbf{استراتژی ماتریسی (Matrix Strategy):} کار \texttt{test} همزمان روی دو مجری مجزا (یکی برای پایتون ۳.۱۰ و دیگری برای ۳.۱۱) اجرا می‌شود.
\item 	extbf{بهینه‌سازی با Caching:} مرحله \texttt{actions/cache} هش فایل \texttt{requirements.txt} را چک می‌کند. اگر لایه‌های پکیج قبلاً دانلود شده باشند، فرایند نصب تنها چند ثانیه زمان می‌برد.
\item 	extbf{همگام‌سازی کارها (\texttt{needs: test}):} کار \texttt{build-and-publish} تا زمانی که کار \texttt{test} با موفقیت کامل تمام نشود، شروع نخواهد شد.
\item 	extbf{کنترل خط‌لوله مبتنی بر شرط (\texttt{if}):} ساخت تصویر داکر فقط زمانی رخ می‌دهد که تغییرات روی شاخه \texttt{main} اعمال شده باشند.
---

\end{enumerate}
\section{خطاهای رایج}

در زمان پیاده‌سازی و مدیریت معماری گیت‌هاب اکشنز، طراحان و مهندسان معمولاً با خطاهایی در زمینه امنیت، کارایی و منطق زیرساخت مواجه می‌شوند.

\subsection{۱. استفاده ناامن از مجری‌های اختصاصی در مخازن عمومی}

یکی از خطرناک‌ترین اشتباهات، اتصال مجری‌های اختصاصی (Self-hosted) به مخازن عمومی (Public Repositories) است. 	extit{ 	extbf{علت:} اگر شخص غریبه‌ای یک \texttt{Pull Request} ایجاد کند، می‌تواند کدهای مخربی در فایل YAML یا تست‌ها قرار دهد که روی ماشین سرور شما (داخل شبکه داخلی سازمان) اجرا شود. } 	extbf{راهکار:} مجری‌های اختصاصی را فقط در مخازن خصوصی (Private) استفاده کنید یا اجرای اتوماتیک اکشن‌ها روی PRهای بیرونی را نیازمند تایید دستی (Manual Approval) کنید.

\subsection{۲. کلیدهای کش نامناسب و عدم انقضا}

عدم تنظیم دقیق کلید کش (\texttt{key}) می‌تواند منجر به استفاده از بسته داده‌های قدیمی (Cache Poisoning یا Stale Cache) شود. 	extit{ 	extbf{علت:} اگر کلید کش شامل هش فایل‌های وابسته (مانند \texttt{package-lock.json} یا \texttt{requirements.txt}) نباشد، مجری همیشه نسخه قدیمی را بازیابی کرده و تغییرات جدید پکیج‌ها اعمال نمی‌شوند. } 	extbf{راهکار:} همواره از توابع اختصاصی مانند \texttt{hashFiles()} در الگوی نام‌گذاری کلید کش استفاده کنید.

\subsection{۳. حلقه‌های بی‌نهایت اجرا (Infinite Loops)}

رخ دادن حالتی که یک اکشن، کدی را commit کرده و به مخزن push می‌کند که دوباره باعث تحریک همان اکشن می‌شود. 	extit{ 	extbf{علت:} عدم کنترل رویداد Push یا عدم استفاده از توکن‌های مناسب. } 	extbf{راهکار:} گیت‌هاب به‌طور خودکار رویدادهای ناشی از \texttt{GITHUB\_TOKEN} را برای جلوگیری از حلقه مسدود می‌کند، اما اگر از یک «توکن دسترسی شخصی» (Personal Access Token - PAT) استفاده کنید، این حلقه رخ خواهد داد. برای جلوگیری، باید پیام‌های commit را فیلتر کنید یا از ابزارهای استاتیک استفاده نمایید.

\subsection{۴. عدم محدودسازی همزمانی (Concurrency Control)}

ارسال چندین Commit متوالی در فاصله زمانی کوتاه می‌تواند ده‌ها مجری را به‌طور همزمان درگیر کند و هزینه‌ها یا منابع سرور را بلعیده و منجر به تداخل در استقرار شود. * 	extbf{راهکار:} استفاده از کلید \texttt{concurrency} در فایل YAML جهت لغو کارهای قدیمی‌ترِ در حال اجرا:

\texttt{}\texttt{yaml concurrency: group: \$\{\{ github.workflow \}\}-\$\{\{ github.ref \}\} cancel-in-progress: true }\texttt{}

---

\section{جمع‌بندی}

گیت‌هاب اکشنز فراتر از یک ابزار CI/CD ساده، یک معماری توزیع‌شده مبتنی بر رویداد است که توانمندی بالایی در اتوماسیون فرایندهای پیچیده مهندسی نرم‌افزار و سیستم‌های هوش مصنوعی ارائه می‌دهد.

در این فصل موارد زیر را فرا گرفتیم: 	extit{ 	extbf{مولفه‌های پایه:} نحوه تعامل گردش‌کارها، کارها، مراحل، اکشن‌ها و مجری‌ها. } 	extbf{معماری سیستم:} تفکیک سطح کنترل و سطح اجرا و مکانیزم ارتباطی Pull بر پایه Long-Polling که امنیت شبکه را تضمین می‌کند. 	extit{ 	extbf{امنیت و کارایی:} نحوه مدیریت رازها با OIDC، تکنیک‌های کش‌سازی و توزیع موازی کارها با استراتژی ماتریسی. } 	extbf{کاربرد عملی:} طراحی یک خط‌لوله کامل شامل تست، تحلیل پوشش کد، و ساخت تصویر داکر.

درک این معماری برای مهندسان سیستم و مدیران زیرساخت ضروری است، چرا که پایه و اساس استقرار خودکار، پایش مدل‌های هوش مصنوعی (MLOps) و مدیریت مقیاس‌پذیر ابری در فصل‌های آتی کتاب خواهد بود.


\end{document}
